% Ajustes: underscores escapados para evitar modo matemático acidental
\chapter{Trabalhos Relacionados}\label{cap_trabalhos_relacionados}

% ------------------------------------------------------------------
% Fluxograma de seleção (adaptado PRISMA 2020)
% Observação: Caso o pacote tikz ainda não esteja no preâmbulo, adicionar:
% \usepackage{tikz}
% \usetikzlibrary{shapes,arrows.meta,positioning}
% (usar arrows.meta garante suporte à ponta 'Latex')
% Se ocorrer erro de pacote ausente ao compilar (ex.: pgf/tikz), instalar via tlmgr:
% tlmgr install pgf
% ------------------------------------------------------------------
% \begin{figure}[H]
% \centering
% \begin{tikzpicture}[
% 	node distance=8mm,
% 	every node/.style={font=\footnotesize},
% 	box/.style={rectangle, draw, rounded corners=2pt, align=center, minimum width=5.4cm, minimum height=0.95cm},
% 	>={Latex[length=3mm]}
% ]

% % Nós
% \node[box] (id) {Identificação\\Registros Google Scholar\\(n = 512)};
% \node[box, below=of id] (triagem) {Triagem (resumos)\\Registros após leitura de resumos\\(n = 23)};
% \node[box, right=24mm of triagem] (excl_resumos) {Excluídos após leitura de resumos\\(n = 489)};

% \node[box, below=14mm of triagem] (texto) {Avaliação em texto completo\\Textos avaliados (n = 23)};
% \node[box, right=24mm of texto] (excl_tema) {Excluídos por desvio temático\\(n = 9)};

% \node[box, below=14mm of texto] (restantes) {Restantes após refino temático\\(n = 14)};
% \node[box, right=24mm of restantes] (excl_app) {Excluídos sem app nativo\\(n = 2)};

% \node[box, below=14mm of restantes] (incluidos) {Trabalhos incluídos na análise comparativa\\(n = 12)};

% % Conexões verticais
% \draw[->] (id) -- (triagem);
% \draw[->] (triagem) -- (texto);
% \draw[->] (texto) -- (restantes);
% \draw[->] (restantes) -- (incluidos);

% % Conexões laterais
% \draw[->] (triagem) -- (excl_resumos);
% \draw[->] (texto) -- (excl_tema);
% \draw[->] (restantes) -- (excl_app);

% \end{tikzpicture}
% \caption{Fluxograma de seleção dos trabalhos (adaptado do modelo PRISMA 2020).}
% \label{fig:fluxo_prisma}
% \end{figure}

O presente capítulo tem como objetivo posicionar o estudo frente às soluções tecnológicas e produções acadêmicas relacionadas, com ênfase em aplicações móveis e sistemas \emph{web} associados. Busca-se identificar padrões, lacunas e oportunidades de aprimoramento, de modo a fundamentar as decisões técnicas e funcionais adotadas no projeto proposto.

O levantamento considerou aplicativos disponíveis nas principais lojas virtuais e trabalhos acadêmicos publicados a partir de 2020, contemplando soluções com temas correlatos ao aleitamento materno e artefatos tecnológicos implementados.

A análise abrangeu dimensões como funcionalidades, suporte profissional, gestão de conteúdo, integrações institucionais e tecnologias empregadas, permitindo estabelecer um panorama comparativo alinhado às demandas atuais de usabilidade, escalabilidade e integração com sistemas públicos de saúde.

\subsection{Levantamento em Lojas de Aplicativos}

Conforme descrito no capítulo de \hyperref[cap_metodologia]{Metodologia}, a primeira abordagem consistiu na realização de uma pesquisa de mercado voltada à identificação de aplicativos móveis relacionados ao aleitamento materno disponíveis nas lojas Google Play e Apple App Store. Essa etapa teve como finalidade mapear soluções já implementadas e compreender como os sistemas subjacentes foram concebidos e desenvolvidos.

\subsubsection{Procedimento de Busca}

A busca foi conduzida utilizando termos e combinações de palavras-chave diretamente associadas ao tema, como "aleitamento materno", "amamentação" e "maternidade". Foram aplicados critérios preliminares de seleção, incluindo: relação direta do aplicativo com o aleitamento materno; existência de indicadores de uso (como número de \emph{downloads} e avaliações); e disponibilidade de um canal de contato com a equipe desenvolvedora, preferencialmente por \emph{e-mail}.

\subsubsection{Tentativa de Contato com Desenvolvedores}

Com o objetivo de aprofundar a compreensão sobre as decisões técnicas e arquiteturais dos aplicativos identificados, foi elaborado um questionário estruturado, disponibilizado via Google Forms em português e inglês. 

O instrumento abordava tópicos como escolhas tecnológicas, arquitetura do sistema, gestão de conteúdo, métricas de uso e desafios enfrentados no desenvolvimento. Aproximadamente 50 equipes foram contatadas, mas não houve retorno, resultando na ausência de dados diretos.

\subsubsection{Implicações da Falta de Respostas}

A ausência de respostas inviabilizou a realização de uma comparação técnica aprofundada, especialmente no que se refere à infraestrutura e às estratégias de atualização das soluções comerciais. Diante dessa limitação, optou-se por reforçar o levantamento documental acadêmico como fonte principal de informações, de modo a suprir as lacunas deixadas pela pesquisa direta junto aos desenvolvedores e garantir a continuidade da análise comparativa.

\subsection{Levantamento Documental}

Para essa etapa, utilizou-se como fonte principal a plataforma Google Scholar. A estratégia de busca envolveu a utilização do texto de pesquisa "aplicativo móvel aleitamento materno web", com filtro temporal a partir de 2020. Como o resultado bruto foi extenso (512 ocorrências), foram aplicados critérios de inclusão e exclusão para refinar a seleção.

\subsubsection{Estratégia de Busca}

Após o resultado da pesquisa inicial, foi realizada uma triagem em múltiplas etapas. A primeira consistiu em percorrer todas as páginas listadas pela fonte, acessar um a um os \emph{links} para os textos acadêmicos, ler seu título, resumo e aplicar os seguintes critérios de inclusão e exclusão:

\begin{itemize}
  \item \textbf{Inclusão:} (a) protótipo ou implementação de aplicativo móvel, com ou sem página \emph{web} associada;
  \item \textbf{Inclusão:} (b) relacionado ao aleitamento materno, maternidade ou saúde da criança;
  \item \textbf{Exclusão:} (c) revisões integrativas puras, ou estudo de casos, sem artefato tecnológico;
  \item \textbf{Exclusão:} (d) temas da área da enfermagem sem relação com tecnologia;
\end{itemize}

Com base na leitura dos resumos, 489 trabalhos foram excluídos por não atenderem aos critérios, restando 23 para avaliação do texto. Nessa segunda etapa, as informações foram extraídas e organizadas em uma planilha, a fim de facilitar a caracterização dos trabalhos. A síntese consolidada é apresentada na Tabela \ref{tab:trabalhos_sintese}.

\begin{small}
\setlength{\tabcolsep}{4pt}
\renewcommand{\arraystretch}{1.06}
\begin{longtable}{p{5.9cm}p{1.9cm}p{1.15cm}p{1.15cm}p{1.25cm}p{1.4cm}p{1.5cm}}
\caption{Síntese dos trabalhos selecionados (2020--2025).}\label{tab:trabalhos_sintese}\\\hline
Título & Tema & App & Painel & Usuários & Suporte & Int. Pública \\\hline
\endfirsthead
\caption[]{(continuação)}\\\hline
Título & Tema & App & Painel & Usuários & Suporte & Int. Pública \\\hline
\endhead
\hline\\[-4pt]
\multicolumn{6}{r}{Continua na próxima página}\\
\endfoot
\hline
\multicolumn{6}{l}{Fonte: Autoria própria (2025).}\\
\endlastfoot
GESTASUS: Aplicativo móvel para integração da caderneta da gestante ao SISPRENATAL WEB & Gestação & Sim & Sim & Sim & Sim & Sim \\
AMAR – Aplicativo de Monitoramento, Acompanhamento e Rastreio do Desenvolvimento Infantil – um estudo de desenvolvimento e validade do conteúdo & Saúde Criança & Sim & Sim & Sim & Sim & Não \\
Avaliação da usabilidade de um aplicativo móvel como facilitador de acesso a serviços de saúde de atenção à gestante de uma maternidade no sul do Maranhão & Gestação & Sim & Sim & Sim & Sim & Não \\
Web app for the monitoring of pregnant and puerperal women: technological production & Gestação & Sim & Sim & Sim & Sim & Não \\
Hora de Mamar! Aplicação Mobile \& Web de acompanhamento alimentar infantil direcionado a pais, cuidadores e profissionais de saúde de bebês de 0 a 12 meses & Aleitamento Materno & Sim & Sim & Não & Sim & Não \\
Sistema AcGest para assistência à saúde de gestantes e puérperas na atenção primária & Gestação & Sim & Sim & Não & Sim & Não \\
Desenvolvimento e validação de aplicativo em saúde no período puerperal & Puerpério & Sim & Não & Sim & Sim & Não \\
Construção de tecnologia m-health para a promoção do aleitamento materno & Aleitamento Materno & Sim & Não & Sim & Sim & Não \\
Aplicativo “Nasci Bem” para profissionais de saúde e familiares de recém-nascidos: construção, validação e avaliação & Recém-nascido & Sim & Não & Sim & Sim & Não \\
Saber G‑estar: construção e validação de um aplicativo móvel para educação em saúde no ciclo gravídico‑puerperal & Gestação & Sim & Não & Sim & Sim & Não \\
Construção e validação de uma tecnologia móvel para apoio às mães de recém-nascidos prematuros & Recém-nascido & Sim & Não & Sim & Sim & Não \\
Desenvolvimento de aplicativo móvel para o acompanhamento pré-natal e validação de conteúdo (Gestação Saudável) & Gestação & Sim & Não & Sim & Sim & Não \\
Produção de aplicativo móvel para orientações a mulher no puerpério & Puerpério & Sim & Não & Não & Sim & Não \\
Uma ontologia de aplicação e plataforma móvel para acompanhamento e promoção do aleitamento materno & Aleitamento Materno & Sim & Não & Não & Sim & Não \\
Desenvolvimento de aplicativo móvel para o acompanhamento pré-natal e validação de conteúdo & Gestação & Sim & Não & Não & Sim & Não \\
Aplicativo móvel sobre a primeira consulta de enfermagem ao recém-nascido na atenção básica: construção e validação & Puericultura & Sim & Não & Não & Sim & Não \\
Aplicativo “PuerpérioSEGURO” em plataforma mobile como tecnologia para o cuidado à beira leito & Puerpério & Sim & Não & Não & Sim & Não \\
SUS\_PRÉ-NATAL\_APP: protótipo de aplicativo móvel para o cuidado pré-natal no Sistema Único de Saúde & Gestação & Sim & Não & Não & Sim & Não \\
PADRÃO OURO: DESENVOLVIMENTO E AVALIAÇÃO DE APLICATIVO MÓVEL SOBRE ALEITAMENTO MATERNO PARA ENFERMEIROS & Aleitamento Materno & Sim & Não & Não & Não & Não \\
SOS MAMA: APLICATIVO MOVEL PARA PUERPERAS QUE VIVENCIAM DIFICULDADES NO ALEITAMENTO MATERNO & Aleitamento Materno & Sim & Não & Não & Não & Não \\
Amamentação com amor: um aplicativo móvel para as lactantes & Aleitamento Materno & Sim & Não & Não & Não & Não \\
Website para família de crianças não amamentadas: desenvolvimento e validação de conteúdo e de interface & Nutrição & Não & Não & Não & Sim & Não \\
\end{longtable}
\end{small}

\vspace{0.4em}
\noindent\textbf{Notas:} "App" indica existência de aplicativo (inclui protótipo); "Suporte" = mediação/atuação de profissional de saúde; "Int. Pública" = integração (prefeituras ou sistemas/sistemas públicos).

Nessa tabela, foi ordenado de cima para baixo por ordem de precedência de "sim" nas colunas Painel, Usuários, Suporte e Integração Pública, de modo a destacar os trabalhos com maior número número de recursos alinhados ao escopo do sistema proposto. Dessa forma, pode-se observar que os 6 primeiros trabalhos são os que mais se aproximam do escopo funcional e técnico desejado, por possuírem painel administrativo, por mais que o tema não seja exclusivamente idêntico ao do presente estudo.

Diante disso, optou-se por selecionar esses 6 trabalhos para uma análise comparativa mais detalhada, conforme apresentado na seção seguinte.

\subsubsection{Caracterização Comparativa}

Esta subseção apresenta a caracterização analítica dos seis trabalhos priorizados, organizando a leitura em três eixos para cada estudo: (i) descrição sintética do propósito e construção; (ii) análise comparativa dos recursos estruturais frente às dimensões funcionais adotadas neste projeto (painel, mediação profissional, gestão de conteúdo e potencial de evolução); e (iii) principais lacunas observadas, agrupadas quando convergem em fragilidades compartilhadas.

\paragraph{GESTASUS}
O GESTASUS é voltado à representação digital da caderneta da gestante, com transferência de dados via QR Code em ambiente autorizado \citeonline{gestasus2025}. Reporta escore SUS médio de 82{,}85 em avaliação com 22 gestantes e 5 especialistas. Fornece ao dispositivo um espelho dos registros clínicos essenciais, acessível localmente, alinhado a fluxos institucionais do pré-natal.

O foco recai na sincronização pontual e visualização estática, sem evidências de curadoria de conteúdo ou versionamento. O vínculo à caderneta reduz necessidade de arquitetura expansível, mas limita flexibilidade modular e interoperabilidade futura.

Ausência de atualização contínua formalizada, ampliação multiprofissional estruturada, trilhas de auditoria e planejamento para adoção intermunicipal.

\paragraph{AMAR}
O sistema AMAR integra módulo \emph{web} para profissionais e aplicativo para famílias, concebido com design participativo \citeonline{amar2021}. Abrange registros clínicos, acompanhamento do desenvolvimento e material multimídia validado. O IVC geral foi $\geq 90\%$, exceto item de diversidade familiar (80\%), com participação de dez profissionais e dez famílias.

Combina registro longitudinal e material educativo, porém sem estratégias declaradas de interoperabilidade ou padrões de integração. Mantém caráter de protótipo, sem processos descritos de evolução sustentada ou validação longitudinal.

Persistem ausência de governança de ciclo de vida do conteúdo e de mecanismos para expansão entre jurisdições ou conexão técnica a sistemas públicos.

\paragraph{Gestação Segura}
O aplicativo para gestantes de maior risco provê suporte comunicacional e acesso segmentado por perfis (profissional e gestante) \citeonline{gestacaosegura2021}. Inclui canal de interação e registros básicos. A usabilidade, avaliada por 20 profissionais, apresentou distribuição SUS: 20\% Bom, 45\% Excelente e 35\% Melhor Alcançável.

Centraliza comunicação e prontuário sintético, com limitação de portabilidade (ausência de múltiplas plataformas) e menor robustez operacional. A camada administrativa concentra-se em inserção assistencial, sem mecanismos claros de ampliação temática ou curadoria de conteúdo estruturada.

Falta de interoperabilidade declarada, plano tecnológico multiambiente e discussão de manutenção preventiva ou observabilidade.

\paragraph{Gestar Care}
O Gestar Care adota abordagem \emph{web} responsiva com áreas usuária, administrativa e clínica \citeonline{gestarcare2022}. A técnica Delphi em duas rodadas produziu concordância: funcionalidade 99\%; confidencialidade/acessibilidade 100\%; viabilidade 85\% (100\% após ajustes); inovação 100\%.

Estrutura multifuncional organizada, porém sem mecanismos descritos de escalonamento ou adaptação regional. Não evidencia padrões de interoperabilidade ou rota de integração incremental nem validação longitudinal.

Ausência de estratégia formal de integração a ecossistemas públicos e de transparência sobre governança de dados sensíveis.

\paragraph{Hora de Mamar}
O Hora de Mamar é aplicação \emph{web} responsiva para registro de eventos alimentares de crianças até doze meses e oferta de central de conhecimento \citeonline{horademamar2021}. Gera relatórios compartilháveis. Não há relato de testes com usuários finais (métrica não reportada).

Predomina lógica transacional de registro e consolidação em relatórios, suportada por conteúdo estático. Não há validação longitudinal em escala nem processo de curadoria contínua explícito.

Falta interoperabilidade, evolução editorial formalizada e evidências de uso monitorado externo ao desenvolvimento.

\paragraph{AcGest}
O AcGest reúne dois aplicativos Android (agente comunitário e gestante/puérpera) e módulo de gestão \emph{web} para consolidação de visitas, ocorrências e cadastros \citeonline{acgest2023}. Foi validado em ambiente simulado com equipe de saúde (sem métricas quantitativas adicionais publicadas).

Múltiplos perfis e coordenação operacional são positivos, porém a dependência de plataforma proprietária aumenta acoplamento e limita portabilidade. Não descreve padrões de interoperabilidade ou intercâmbio de dados.

Falta interoperabilidade formal, validação em escala territorial e há risco adicional de dependência tecnológica concentrada.

\paragraph{Síntese das Fragilidades Comuns}
Observam-se padrões transversais: (i) inexistência ou descrição incipiente de interoperabilidade com sistemas públicos; (ii) curadoria de conteúdo e versionamento contínuo pouco formalizados; (iii) validações pontuais sem validação longitudinal; e (iv) escassez de estratégias técnicas para evolução modular e adoção intermunicipal. Esse agrupamento sintetiza o racional comparativo sem reiterar individualmente cada fragilidade já compartilhada.
